% !TEX root = ../main.tex
\section{Reproduction results}
\label{sec:reproduction}

We run $N \in \{1,2,4,8\}$ on Swimmer-v4, Hopper-v4 and Humanoid-v4 with five
random seeds per configuration (60 runs). Table~\ref{tab:reproduction} reports
the mean of each seed's final 100 iterations, then the mean and standard
deviation across seeds, the same quantity Table~2 of the original reports.

\begin{table}[t]
  \centering
  \caption{Final rewards, mean $\pm$ standard deviation across five seeds.
  Original values are from Table~2 of \citet{lan2023fednpgadmm}, computed over
  ten seeds. Humanoid does not appear in their Table~2.}
  \label{tab:reproduction}
  \begin{tabular}{llrrrr}
    \toprule
    Environment & Source & $N=1$ & $N=2$ & $N=4$ & $N=8$ \\
    \midrule
    \multirow{2}{*}{Swimmer-v4}
      & Ours  & $\mathbf{123.9 \pm 2.3}$ & $126.0 \pm 1.3$ & $127.0 \pm 0.8$ & $\mathbf{128.4 \pm 0.8}$ \\
      & Paper & $109.4 \pm 5.9$  & $123.6 \pm 10.3$ & $127.2 \pm 2.2$ & $128.5 \pm 1.9$ \\
    \midrule
    \multirow{2}{*}{Hopper-v4}
      & Ours  & $\mathbf{2497 \pm 635}$ & $2140 \pm 386$ & $1658 \pm 273$ & $1605 \pm 159$ \\
      & Paper & $1473 \pm 383$ & $2384 \pm 371$ & $2507 \pm 230$ & $2719 \pm 173$ \\
    \midrule
    Humanoid-v4
      & Ours  & $993 \pm 187$ & $3391 \pm 796$ & $5343 \pm 121$ & $\mathbf{6329 \pm 334}$ \\
    \bottomrule
  \end{tabular}
\end{table}

\begin{figure}[t]
  \centering
  \begin{subfigure}[b]{0.49\textwidth}
    %% Convergence, Swimmer-v4. Mean over five seeds; band spans min to max.
%% Data: figures/convergence_swimmer.dat (scripts/export_figures.py).
\begin{tikzpicture}
\begin{axis}[
    width=\linewidth,
    height=0.78\linewidth,
    xlabel={Iteration},
    ylabel={Average batch return},
    xmin=1, xmax=1000,
    ymin=0, ymax=140,
    xtick distance=250,
    ytick distance=35,
    ymajorgrids=true, xmajorgrids=true, grid style=dashed,
    tick label style={font=\footnotesize},
    label style={font=\footnotesize},
    legend style={legend columns=2, font=\footnotesize, cells={anchor=west},
                  anchor=south east, at={(0.99,0.02)}},
]

\addplot[name path=N1lo, draw=none, forget plot]
    table[x=iteration, y=N1_lo] {figures/convergence_swimmer.dat};
\addplot[name path=N1hi, draw=none, forget plot]
    table[x=iteration, y=N1_hi] {figures/convergence_swimmer.dat};
\addplot[colorA, opacity=0.15, forget plot] fill between[of=N1lo and N1hi];
\addplot[colorA, line width=0.9pt]
    table[x=iteration, y=N1_mean] {figures/convergence_swimmer.dat};
\addlegendentry{$N=1$};

\addplot[name path=N2lo, draw=none, forget plot]
    table[x=iteration, y=N2_lo] {figures/convergence_swimmer.dat};
\addplot[name path=N2hi, draw=none, forget plot]
    table[x=iteration, y=N2_hi] {figures/convergence_swimmer.dat};
\addplot[colorB, opacity=0.15, forget plot] fill between[of=N2lo and N2hi];
\addplot[colorB, line width=0.9pt]
    table[x=iteration, y=N2_mean] {figures/convergence_swimmer.dat};
\addlegendentry{$N=2$};

\addplot[name path=N4lo, draw=none, forget plot]
    table[x=iteration, y=N4_lo] {figures/convergence_swimmer.dat};
\addplot[name path=N4hi, draw=none, forget plot]
    table[x=iteration, y=N4_hi] {figures/convergence_swimmer.dat};
\addplot[colorC, opacity=0.15, forget plot] fill between[of=N4lo and N4hi];
\addplot[colorC, line width=0.9pt]
    table[x=iteration, y=N4_mean] {figures/convergence_swimmer.dat};
\addlegendentry{$N=4$};

\addplot[name path=N8lo, draw=none, forget plot]
    table[x=iteration, y=N8_lo] {figures/convergence_swimmer.dat};
\addplot[name path=N8hi, draw=none, forget plot]
    table[x=iteration, y=N8_hi] {figures/convergence_swimmer.dat};
\addplot[colorD, opacity=0.15, forget plot] fill between[of=N8lo and N8hi];
\addplot[colorD, line width=0.9pt]
    table[x=iteration, y=N8_mean] {figures/convergence_swimmer.dat};
\addlegendentry{$N=8$};
\end{axis}
\end{tikzpicture}

    \caption{Swimmer-v4}
  \end{subfigure}
  \hfill
  \begin{subfigure}[b]{0.49\textwidth}
    %% Convergence, Hopper-v4. Mean over five seeds; band spans min to max.
%% Data: figures/convergence_hopper.dat (scripts/export_figures.py).
\begin{tikzpicture}
\begin{axis}[
    width=\linewidth,
    height=0.78\linewidth,
    xlabel={Iteration},
    ylabel={Average batch return},
    xmin=1, xmax=2000,
    ymin=0, ymax=3600,
    xtick distance=500,
    ytick distance=900,
    scaled y ticks=base 10:-3,
    ymajorgrids=true, xmajorgrids=true, grid style=dashed,
    tick label style={font=\footnotesize},
    label style={font=\footnotesize},
    legend style={legend columns=2, font=\footnotesize, cells={anchor=west},
                  anchor=south east, at={(0.99,0.02)}},
]

\addplot[name path=N1lo, draw=none, forget plot]
    table[x=iteration, y=N1_lo] {figures/convergence_hopper.dat};
\addplot[name path=N1hi, draw=none, forget plot]
    table[x=iteration, y=N1_hi] {figures/convergence_hopper.dat};
\addplot[colorA, opacity=0.15, forget plot] fill between[of=N1lo and N1hi];
\addplot[colorA, line width=0.9pt]
    table[x=iteration, y=N1_mean] {figures/convergence_hopper.dat};
\addlegendentry{$N=1$};

\addplot[name path=N2lo, draw=none, forget plot]
    table[x=iteration, y=N2_lo] {figures/convergence_hopper.dat};
\addplot[name path=N2hi, draw=none, forget plot]
    table[x=iteration, y=N2_hi] {figures/convergence_hopper.dat};
\addplot[colorB, opacity=0.15, forget plot] fill between[of=N2lo and N2hi];
\addplot[colorB, line width=0.9pt]
    table[x=iteration, y=N2_mean] {figures/convergence_hopper.dat};
\addlegendentry{$N=2$};

\addplot[name path=N4lo, draw=none, forget plot]
    table[x=iteration, y=N4_lo] {figures/convergence_hopper.dat};
\addplot[name path=N4hi, draw=none, forget plot]
    table[x=iteration, y=N4_hi] {figures/convergence_hopper.dat};
\addplot[colorC, opacity=0.15, forget plot] fill between[of=N4lo and N4hi];
\addplot[colorC, line width=0.9pt]
    table[x=iteration, y=N4_mean] {figures/convergence_hopper.dat};
\addlegendentry{$N=4$};

\addplot[name path=N8lo, draw=none, forget plot]
    table[x=iteration, y=N8_lo] {figures/convergence_hopper.dat};
\addplot[name path=N8hi, draw=none, forget plot]
    table[x=iteration, y=N8_hi] {figures/convergence_hopper.dat};
\addplot[colorD, opacity=0.15, forget plot] fill between[of=N8lo and N8hi];
\addplot[colorD, line width=0.9pt]
    table[x=iteration, y=N8_mean] {figures/convergence_hopper.dat};
\addlegendentry{$N=8$};
\end{axis}
\end{tikzpicture}

    \caption{Hopper-v4}
  \end{subfigure}

  \vspace{0.6em}
  \begin{subfigure}[b]{0.49\textwidth}
    %% Convergence, Humanoid-v4. Mean over five seeds; band spans min to max.
%% Data: figures/convergence_humanoid.dat (scripts/export_figures.py).
\begin{tikzpicture}
\begin{axis}[
    width=\linewidth,
    height=0.78\linewidth,
    xlabel={Iteration},
    ylabel={Average batch return},
    xmin=1, xmax=3000,
    ymin=0, ymax=7000,
    xtick distance=1000,
    ytick distance=1750,
    scaled y ticks=base 10:-3,
    ymajorgrids=true, xmajorgrids=true, grid style=dashed,
    tick label style={font=\footnotesize},
    label style={font=\footnotesize},
    legend style={legend columns=2, font=\footnotesize, cells={anchor=west},
                  anchor=north west, at={(0.02,0.98)}},
]

\addplot[name path=N1lo, draw=none, forget plot]
    table[x=iteration, y=N1_lo] {figures/convergence_humanoid.dat};
\addplot[name path=N1hi, draw=none, forget plot]
    table[x=iteration, y=N1_hi] {figures/convergence_humanoid.dat};
\addplot[colorA, opacity=0.15, forget plot] fill between[of=N1lo and N1hi];
\addplot[colorA, line width=0.9pt]
    table[x=iteration, y=N1_mean] {figures/convergence_humanoid.dat};
\addlegendentry{$N=1$};

\addplot[name path=N2lo, draw=none, forget plot]
    table[x=iteration, y=N2_lo] {figures/convergence_humanoid.dat};
\addplot[name path=N2hi, draw=none, forget plot]
    table[x=iteration, y=N2_hi] {figures/convergence_humanoid.dat};
\addplot[colorB, opacity=0.15, forget plot] fill between[of=N2lo and N2hi];
\addplot[colorB, line width=0.9pt]
    table[x=iteration, y=N2_mean] {figures/convergence_humanoid.dat};
\addlegendentry{$N=2$};

\addplot[name path=N4lo, draw=none, forget plot]
    table[x=iteration, y=N4_lo] {figures/convergence_humanoid.dat};
\addplot[name path=N4hi, draw=none, forget plot]
    table[x=iteration, y=N4_hi] {figures/convergence_humanoid.dat};
\addplot[colorC, opacity=0.15, forget plot] fill between[of=N4lo and N4hi];
\addplot[colorC, line width=0.9pt]
    table[x=iteration, y=N4_mean] {figures/convergence_humanoid.dat};
\addlegendentry{$N=4$};

\addplot[name path=N8lo, draw=none, forget plot]
    table[x=iteration, y=N8_lo] {figures/convergence_humanoid.dat};
\addplot[name path=N8hi, draw=none, forget plot]
    table[x=iteration, y=N8_hi] {figures/convergence_humanoid.dat};
\addplot[colorD, opacity=0.15, forget plot] fill between[of=N8lo and N8hi];
\addplot[colorD, line width=0.9pt]
    table[x=iteration, y=N8_mean] {figures/convergence_humanoid.dat};
\addlegendentry{$N=8$};
\end{axis}
\end{tikzpicture}

    \caption{Humanoid-v4}
  \end{subfigure}
  \caption{Convergence for each agent count. The solid line is the mean over five
  seeds and the shaded band spans the minimum and maximum seed, so the band
  shows the full observed range and not a modelled interval. Curves are
  smoothed with a centred 21-iteration rolling mean before aggregation,
  following \citet{spinningup}, who smooth the same quantity over eleven
  epochs. Swimmer and Humanoid separate cleanly by agent count. On Hopper the
  bands overlap throughout and the ordering reverses, which is the visual form
  of the result discussed in Section~\ref{sec:hopper}.}
  \label{fig:convergence}
\end{figure}

\begin{figure}[t]
  \centering
  \begin{subfigure}[b]{0.49\textwidth}
    %% Return against cumulative uplink traffic, Swimmer-v4.
%% x is the federation total: the per-agent cost is 2d*8 bytes per iteration
%% for every N, and that invariance is the O(d) claim itself.
%% Data: figures/communication_swimmer.dat (scripts/export_figures.py).
\begin{tikzpicture}
\begin{axis}[
    width=\linewidth,
    height=0.78\linewidth,
    xlabel={Cumulative uplink (bytes)},
    ylabel={Average batch return},
    xmode=log,
    ymin=0, ymax=140,
    ytick distance=35,
    ymajorgrids=true, xmajorgrids=true, grid style=dashed,
    tick label style={font=\footnotesize},
    label style={font=\footnotesize},
    legend style={legend columns=2, font=\footnotesize, cells={anchor=west},
                  anchor=north west, at={(0.02,0.98)}},
]

\addplot[name path=N1clo, draw=none, forget plot]
    table[x=N1_bytes, y=N1_lo] {figures/communication_swimmer.dat};
\addplot[name path=N1chi, draw=none, forget plot]
    table[x=N1_bytes, y=N1_hi] {figures/communication_swimmer.dat};
\addplot[colorA, opacity=0.15, forget plot] fill between[of=N1clo and N1chi];
\addplot[colorA, line width=0.9pt]
    table[x=N1_bytes, y=N1_mean] {figures/communication_swimmer.dat};
\addlegendentry{$N=1$};

\addplot[name path=N2clo, draw=none, forget plot]
    table[x=N2_bytes, y=N2_lo] {figures/communication_swimmer.dat};
\addplot[name path=N2chi, draw=none, forget plot]
    table[x=N2_bytes, y=N2_hi] {figures/communication_swimmer.dat};
\addplot[colorB, opacity=0.15, forget plot] fill between[of=N2clo and N2chi];
\addplot[colorB, line width=0.9pt]
    table[x=N2_bytes, y=N2_mean] {figures/communication_swimmer.dat};
\addlegendentry{$N=2$};

\addplot[name path=N4clo, draw=none, forget plot]
    table[x=N4_bytes, y=N4_lo] {figures/communication_swimmer.dat};
\addplot[name path=N4chi, draw=none, forget plot]
    table[x=N4_bytes, y=N4_hi] {figures/communication_swimmer.dat};
\addplot[colorC, opacity=0.15, forget plot] fill between[of=N4clo and N4chi];
\addplot[colorC, line width=0.9pt]
    table[x=N4_bytes, y=N4_mean] {figures/communication_swimmer.dat};
\addlegendentry{$N=4$};

\addplot[name path=N8clo, draw=none, forget plot]
    table[x=N8_bytes, y=N8_lo] {figures/communication_swimmer.dat};
\addplot[name path=N8chi, draw=none, forget plot]
    table[x=N8_bytes, y=N8_hi] {figures/communication_swimmer.dat};
\addplot[colorD, opacity=0.15, forget plot] fill between[of=N8clo and N8chi];
\addplot[colorD, line width=0.9pt]
    table[x=N8_bytes, y=N8_mean] {figures/communication_swimmer.dat};
\addlegendentry{$N=8$};
\end{axis}
\end{tikzpicture}

    \caption{Swimmer-v4}
  \end{subfigure}
  \hfill
  \begin{subfigure}[b]{0.49\textwidth}
    %% Return against cumulative uplink traffic, Hopper-v4.
%% x is the federation total: the per-agent cost is 2d*8 bytes per iteration
%% for every N, and that invariance is the O(d) claim itself.
%% Data: figures/communication_hopper.dat (scripts/export_figures.py).
\begin{tikzpicture}
\begin{axis}[
    width=\linewidth,
    height=0.78\linewidth,
    xlabel={Cumulative uplink (bytes)},
    ylabel={Average batch return},
    xmode=log,
    ymin=0, ymax=3600,
    ytick distance=900,
    scaled y ticks=base 10:-3,
    ymajorgrids=true, xmajorgrids=true, grid style=dashed,
    tick label style={font=\footnotesize},
    label style={font=\footnotesize},
    legend style={legend columns=2, font=\footnotesize, cells={anchor=west},
                  anchor=north west, at={(0.02,0.98)}},
]

\addplot[name path=N1clo, draw=none, forget plot]
    table[x=N1_bytes, y=N1_lo] {figures/communication_hopper.dat};
\addplot[name path=N1chi, draw=none, forget plot]
    table[x=N1_bytes, y=N1_hi] {figures/communication_hopper.dat};
\addplot[colorA, opacity=0.15, forget plot] fill between[of=N1clo and N1chi];
\addplot[colorA, line width=0.9pt]
    table[x=N1_bytes, y=N1_mean] {figures/communication_hopper.dat};
\addlegendentry{$N=1$};

\addplot[name path=N2clo, draw=none, forget plot]
    table[x=N2_bytes, y=N2_lo] {figures/communication_hopper.dat};
\addplot[name path=N2chi, draw=none, forget plot]
    table[x=N2_bytes, y=N2_hi] {figures/communication_hopper.dat};
\addplot[colorB, opacity=0.15, forget plot] fill between[of=N2clo and N2chi];
\addplot[colorB, line width=0.9pt]
    table[x=N2_bytes, y=N2_mean] {figures/communication_hopper.dat};
\addlegendentry{$N=2$};

\addplot[name path=N4clo, draw=none, forget plot]
    table[x=N4_bytes, y=N4_lo] {figures/communication_hopper.dat};
\addplot[name path=N4chi, draw=none, forget plot]
    table[x=N4_bytes, y=N4_hi] {figures/communication_hopper.dat};
\addplot[colorC, opacity=0.15, forget plot] fill between[of=N4clo and N4chi];
\addplot[colorC, line width=0.9pt]
    table[x=N4_bytes, y=N4_mean] {figures/communication_hopper.dat};
\addlegendentry{$N=4$};

\addplot[name path=N8clo, draw=none, forget plot]
    table[x=N8_bytes, y=N8_lo] {figures/communication_hopper.dat};
\addplot[name path=N8chi, draw=none, forget plot]
    table[x=N8_bytes, y=N8_hi] {figures/communication_hopper.dat};
\addplot[colorD, opacity=0.15, forget plot] fill between[of=N8clo and N8chi];
\addplot[colorD, line width=0.9pt]
    table[x=N8_bytes, y=N8_mean] {figures/communication_hopper.dat};
\addlegendentry{$N=8$};
\end{axis}
\end{tikzpicture}

    \caption{Hopper-v4}
  \end{subfigure}

  \vspace{0.6em}
  \begin{subfigure}[b]{0.49\textwidth}
    %% Return against cumulative uplink traffic, Humanoid-v4.
%% x is the federation total: the per-agent cost is 2d*8 bytes per iteration
%% for every N, and that invariance is the O(d) claim itself.
%% Data: figures/communication_humanoid.dat (scripts/export_figures.py).
\begin{tikzpicture}
\begin{axis}[
    width=\linewidth,
    height=0.78\linewidth,
    xlabel={Cumulative uplink (bytes)},
    ylabel={Average batch return},
    xmode=log,
    ymin=0, ymax=7000,
    ytick distance=1750,
    scaled y ticks=base 10:-3,
    ymajorgrids=true, xmajorgrids=true, grid style=dashed,
    tick label style={font=\footnotesize},
    label style={font=\footnotesize},
    legend style={legend columns=2, font=\footnotesize, cells={anchor=west},
                  anchor=north west, at={(0.02,0.98)}},
]

\addplot[name path=N1clo, draw=none, forget plot]
    table[x=N1_bytes, y=N1_lo] {figures/communication_humanoid.dat};
\addplot[name path=N1chi, draw=none, forget plot]
    table[x=N1_bytes, y=N1_hi] {figures/communication_humanoid.dat};
\addplot[colorA, opacity=0.15, forget plot] fill between[of=N1clo and N1chi];
\addplot[colorA, line width=0.9pt]
    table[x=N1_bytes, y=N1_mean] {figures/communication_humanoid.dat};
\addlegendentry{$N=1$};

\addplot[name path=N2clo, draw=none, forget plot]
    table[x=N2_bytes, y=N2_lo] {figures/communication_humanoid.dat};
\addplot[name path=N2chi, draw=none, forget plot]
    table[x=N2_bytes, y=N2_hi] {figures/communication_humanoid.dat};
\addplot[colorB, opacity=0.15, forget plot] fill between[of=N2clo and N2chi];
\addplot[colorB, line width=0.9pt]
    table[x=N2_bytes, y=N2_mean] {figures/communication_humanoid.dat};
\addlegendentry{$N=2$};

\addplot[name path=N4clo, draw=none, forget plot]
    table[x=N4_bytes, y=N4_lo] {figures/communication_humanoid.dat};
\addplot[name path=N4chi, draw=none, forget plot]
    table[x=N4_bytes, y=N4_hi] {figures/communication_humanoid.dat};
\addplot[colorC, opacity=0.15, forget plot] fill between[of=N4clo and N4chi];
\addplot[colorC, line width=0.9pt]
    table[x=N4_bytes, y=N4_mean] {figures/communication_humanoid.dat};
\addlegendentry{$N=4$};

\addplot[name path=N8clo, draw=none, forget plot]
    table[x=N8_bytes, y=N8_lo] {figures/communication_humanoid.dat};
\addplot[name path=N8chi, draw=none, forget plot]
    table[x=N8_bytes, y=N8_hi] {figures/communication_humanoid.dat};
\addplot[colorD, opacity=0.15, forget plot] fill between[of=N8clo and N8chi];
\addplot[colorD, line width=0.9pt]
    table[x=N8_bytes, y=N8_mean] {figures/communication_humanoid.dat};
\addlegendentry{$N=8$};
\end{axis}
\end{tikzpicture}

    \caption{Humanoid-v4}
  \end{subfigure}
  \caption{The same runs plotted against cumulative uplink traffic instead of
  iteration count, on a logarithmic axis. The horizontal axis is the total
  across the federation: each agent transmits $2d$ values per iteration
  whatever $N$ is, so a per-agent axis would place all four series on top of
  one another, and that invariance is the $\Oh(d)$ property itself. Read this
  way the figure prices the extra agents, and the curves shift right as $N$
  grows.}
  \label{fig:communication}
\end{figure}

\subsection{Swimmer: reproduced}

Figure~\ref{fig:convergence}(a) shows the curves behind the Swimmer column, and
Figure~\ref{fig:communication}(a) the same runs against uplink traffic.
Our numbers fall within the reported confidence intervals at every agent
count, they are monotone in $N$ as the paper describes, and their seed variance
is in fact \emph{tighter} than the original's at every count, by a factor of
roughly two at both ends ($\pm 2.3$ against $\pm 5.9$ at $N=1$, $\pm 0.8$
against $\pm 1.9$ at $N=8$). At $N=1$ we also exceed the reported mean, $123.9$
against $109.4$. On this environment we therefore regard the result as
reproduced without qualification.

The variance does still fall with $N$, from $\pm 2.3$ at one agent to $\pm 0.8$
at eight, which is the direction the paper reports. The effect is modest here
because every seed converges: none of the five single-agent runs falls into the
local optimum Swimmer is known for, so there is no failure mode for the extra
agents to suppress. Humanoid in Section~\ref{sec:humanoid} shows the same claim
with far more room to act.

\subsection{Humanoid: strong scaling}
\label{sec:humanoid}

Humanoid is absent from the original Table~2, so we have no direct comparison.
The qualitative claim nevertheless holds, since return grows $6.4\times$ from $N=1$ to $N=8$,
monotonically, with the $N=4$ and $N=8$ variance small relative to the mean.

\subsection{Hopper: not reproduced}
\label{sec:hopper}

Hopper is the one environment where our result contradicts the original. Our
$N=1$ \emph{exceeds} the reported value ($2497$ against $1473$), but
performance then \emph{decreases} monotonically with more agents, ending at
$1605$ where the paper reports $2719$. The trend is inverted.

Section~\ref{sec:discussion} lists the aggregation hypotheses we tested and
refuted. The explanation consistent with the evidence is a
speed-versus-survival trade-off specific to Hopper's reward structure.
Table~\ref{tab:hopper} shows that the runs with more agents earn \emph{more}
reward per timestep while surviving \emph{fewer} timesteps: they converge to a
faster, less stable gait.

\begin{table}[t]
  \centering
  \caption{Hopper-v4: reward per step against episode length, final 100
  iterations. More agents yield a faster gait that falls sooner.}
  \label{tab:hopper}
  \begin{tabular}{lrrr}
    \toprule
    Configuration & Reward/step & Episode length & Return \\
    \midrule
    $N=1$ (best seed) & $3.76$ & $779$ & $2928$ \\
    $N=2$ & $4.20$ & $372$ & $1562$ \\
    $N=4$ & $4.27$ & $384$ & $1638$ \\
    $N=8$ & $\mathbf{4.49}$ & $\mathbf{368}$ & $1652$ \\
    \bottomrule
  \end{tabular}
\end{table}

Hopper's \texttt{healthy\_reward} is $1.0$ per surviving step, small relative
to its forward-velocity reward, so a gait that runs fast and falls at step
$\sim$400 is a local optimum. Humanoid's \texttt{healthy\_reward} is $5.0$,
which aligns survival with scoring, and there reward per step and episode
length rise
\emph{together} with $N$ ($5.70 \to 7.09$ and $206 \to 874$). The mechanism we propose is that additional agents reduce gradient noise, and
with it the implicit exploration that noise supplies. Whether the loss of that
exploration helps or hurts then depends on which basin the now more decisive
optimizer settles into. On Swimmer and Humanoid it steers away from the poor
optimum, and the seed variance falls accordingly; on Hopper it steers into one.
The same property of the method therefore produces the paper's result on two
environments and the opposite on the third.

\subsection{On the variance of the reported metric}
\label{sec:variance}

Single-seed Hopper and Humanoid curves swing between roughly $1000$ and $3700$
from one iteration to the next, while Swimmer's are flat to within $2\%$. This
is not instability of the algorithm. Swimmer-v4 has no termination condition
at all, and \citet{gymnasium} state that ``the Swimmer never terminates'', so
every episode runs exactly $1000$ steps. Hopper and Humanoid terminate when
the robot falls, which makes return approximately affine in survival time and
the per-batch mean a draw from a two-component mixture: roughly $58\%$ of
episodes reach the time limit near $3700$, while $9\%$ fall early below
$2000$. With one to two episodes per plotted point, the swing is the expected
behaviour of the estimator. The same policy evaluated deterministically scores
$3749 \pm 4$ and almost never falls.

We verified that the swing cannot be tuned away. Removing the
exploration-noise floor, removing the ceiling, and annealing the noise to a
level at which a fixed policy provably never falls all left the fall rate at
$9$--$12\%$. A policy trained at lower noise re-optimizes to a correspondingly
more marginal gait, so the failure rate is an equilibrium of the objective
rather than a function of the noise level. A never-falling policy would score
lower in expectation.
